\section{The 45-Minute Market Discovery Protocol}

Anchor every customer discovery interview upon a concrete historical event, rather than asking prospective customers to speculate about their hypothetical future desires:

\begin{enumerate}[leftmargin=14pt,itemsep=3pt,topsep=2pt]
  \item \textbf{The Anchor Event (Minutes 1 to 5):} ``Think back to the most recent time this specific technical problem became painful enough that you actively sought a solution. What exact event triggered that search last month?''
  \item \textbf{Consequences and Inaction (Minutes 6 to 15):} ``What were the operational stakes if the problem remained unresolved? What did you attempt first? What would have happened to your team if you simply did nothing?''
  \item \textbf{The Procurement Mechanics (Minutes 16 to 25):} ``Who had to approve the expenditure? Did the funds come from personal savings, departmental expense accounts, or corporate HR? What security review, discount demand, or committee vote almost killed the transaction?''
  \item \textbf{Dissatisfaction and Workarounds (Minutes 26 to 35):} ``What remained broken or unaddressed after using that solution? What manual scripts, extra hours, or awkward workarounds does your team still endure today?''
  \item \textbf{The Competitive Ecosystem (Minutes 36 to 40):} ``Who else in your professional network experiences this exact operational bottleneck? Which vendor or authority do you genuinely respect in this domain, and why?''
  \item \textbf{The Diagnostic Close (Minutes 41 to 45):} ``What critical operational nuance did I fail to ask about today? If I design a tightly bounded, paid technical diagnostic clinic specifically addressing this failure mode, may I share the syllabus with you for an objective critique?''
\end{enumerate}

\textbf{Never pitch your product during a discovery interview.} If you disguise a sales presentation as research, customers will politely flatter your ego while providing utterly useless behavioral data.

\section{The 30-Day Market Investigation Campaign}

A founder operating with strictly limited weekly hours can complete a comprehensive market investigation within a single 30-day calendar sprint:

\begin{itemize}[leftmargin=12pt,itemsep=3pt,topsep=2pt]
  \item \textbf{Days 1 to 5 (Hypotheses \& Boundary Filters):} Draft 4 to 6 distinct buyer-job-trigger hypotheses. Pass them through the 6 non-negotiable boundary gates. Discard all options requiring heavy upfront capital, regulatory licensing, or excessive travel.
  \item \textbf{Days 6 to 12 (Base-Rate Research \& Prospect Mapping):} Extract verified public statistics, competitor pricing schedules, and active job postings. Compile a named list of 40 qualified prospective decision-makers.
  \item \textbf{Days 13 to 20 (Qualitative Field Interviews):} Execute 12 in-depth discovery interviews with practicing engineers, academic chairs, and engineering directors. Reconstruct their exact spending history and active manual workarounds.
  \item \textbf{Days 21 to 26 (Commercial Pilot Offer Design):} Package a minimal, complete commercial offering (an 8-seat weekend diagnostic clinic or a single-semester lab proposal). Model direct variable costs, TA expenses, and founder hours.
  \item \textbf{Days 27 to 30 (Executive Decision Memo):} Synthesize verified behavioral data. If qualified buyers commit funds, proceed to fulfillment. If buyers refuse to pay, pivot to the next candidate market immediately.
\end{itemize}

\section{The Opportunity-Selection Decision Tree}

Evaluate every candidate commercial opportunity through this disciplined sequential logic tree:

\begin{enumerate}[leftmargin=14pt,itemsep=3pt,topsep=2pt]
  \item \textbf{Step 1: Can you identify a concrete decision-maker and an acute, consequential operational job?}
  \begin{itemize}[leftmargin=10pt,itemsep=1pt]
    \item \textit{If No:} Halt. Continue qualitative field discovery before writing software.
    \item \textit{If Yes:} Proceed to Step 2.
  \end{itemize}

  \item \textbf{Step 2: Is there verifiable empirical evidence of current capital expenditure, expensive manual workarounds, or painful friction?}
  \begin{itemize}[leftmargin=10pt,itemsep=1pt]
    \item \textit{If No:} Halt. You have discovered a mild hobby, not a commercial market.
    \item \textit{If Yes:} Proceed to Step 3.
  \end{itemize}

  \item \textbf{Step 3: Can the commercial transaction execute cleanly within existing legal, budgetary, and institutional constraints?}
  \begin{itemize}[leftmargin=10pt,itemsep=1pt]
    \item \textit{If No:} Alter the economic buyer, shift to an adjacent value-chain layer, or defer.
    \item \textit{If Yes:} Proceed to Step 4.
  \end{itemize}

  \item \textbf{Step 4: Can the founding team economically reach enough qualified buyers through verified channels?}
  \begin{itemize}[leftmargin=10pt,itemsep=1pt]
    \item \textit{If No:} Narrow the initial customer niche or partner with an entity controlling distribution.
    \item \textit{If Yes:} Proceed to Step 5.
  \end{itemize}

  \item \textbf{Step 5: Can you deliver the promised transformation within founder weekly hours while retaining healthy contribution margin?}
  \begin{itemize}[leftmargin=10pt,itemsep=1pt]
    \item \textit{If No:} Redesign delivery protocols, automate evaluation rubrics, or raise prices.
    \item \textit{If Yes:} Proceed to Step 6.
  \end{itemize}

  \item \textbf{Step 6: Does the opportunity survive rational incumbent reactions and build reusable strategic assets for the enduring institution?}
  \begin{itemize}[leftmargin=10pt,itemsep=1pt]
    \item \textit{If No:} Re-examine positioning or select a more defensible niche.
    \item \textit{If Yes:} \textbf{Commit capital to the 30-day paid commercial experiment immediately.}
  \end{itemize}
\end{enumerate}
